\begin{proof}[Proof of \cref{obj:lem:public-error-certificate}]
\begin{enumerate}
\item
Index the observed records and form the three blocks as in
\cref{obj:synth_3,obj:def:total-estimator}. Thus
\[
O_i=(X_i,W_i,Y_i),\qquad V_i=W_i-a_0,\qquad a_0=\frac12,
\qquad i=0,\ldots,n-1,
\]
and write
\[
N_r:=|\mathcal I_r|,
\qquad
\mathcal I_r=\{i\in\{0,\ldots,n-1\}:i\bmod3=r\},
\qquad r\in\{0,1,2\}.
\]
The indices \(3i+r\), for \(0\le i<\lfloor n/3\rfloor\), belong to
\(\mathcal I_r\). Consequently,
\[
N_r\ge\left\lfloor\frac n3\right\rfloor\ge1,
\qquad
n\le3\left\lfloor\frac n3\right\rfloor+2
\le5\left\lfloor\frac n3\right\rfloor
\le5N_r.
\]
In particular, every block is nonempty. The positive weight moment and public constants give
\[
b_q(K)=p_{\min}c_{\mathrm{lo}}I_\kappa(q)>0.
\]
For each \(x\in\mathcal X\), define the population moments
\[
D_x(q):=E_P[\mathbf1\{X=x\}\ell_q(V)],
\qquad
M_x(q):=E_P[\mathbf1\{X=x\}Y\ell_q(V)].
\]
The empirical moments are
\[
\widehat p_x=\frac1{N_0}\sum_{i\in\mathcal I_0}\mathbf1\{X_i=x\},
\]
\[
\widehat D_x(q)=\frac1{N_1}\sum_{i\in\mathcal I_1}
\mathbf1\{X_i=x\}\ell_q(V_i),
\qquad
\widehat M_x(q)=\frac1{N_2}\sum_{i\in\mathcal I_2}
\mathbf1\{X_i=x\}Y_i\ell_q(V_i).
\]
On every observed sample in \((\mathcal X\times\mathbb R\times\mathbb R)^n\), define
\[
\Delta_{D,x}:=|\widehat D_x(q)-D_x(q)|,\qquad
\Delta_{M,x}:=|\widehat M_x(q)-M_x(q)|,\qquad
\Delta_{p,x}:=|\widehat p_x-p_x|.
\]
All sample expectations below are taken under \(Q_P^n\).
% lean: CausalSmith.Stat.NoisydoseWeakdesignTransition.blocksNonempty_of_three_le
% lean: CausalSmith.Stat.NoisydoseWeakdesignTransition.splitBlock_card_lower
% lean: CausalSmith.Stat.NoisydoseWeakdesignTransition.splitBlock_card_fifth

\item
We first establish the population moments of these averages. Choose the
conditional densities \(g_x\) supplied by \cref{obj:ass:weak-design}.
Define the standard Gaussian probability measure by
\[
\Phi_1:=\mathcal N(0,1).
\]
The Gaussian channel and error independence imply that, conditional on
\(X=x\), the joint law of \(Z\) and \(A-a_0\) is the product of
\(\Phi_1\) and the measure with density \(g_x\) on \([-1/2,1/2]\).
The public second-moment quantity in \cref{obj:synth_4} satisfies
\[
V_q(K)
=c_{\mathrm{hi}}\int_{-1/2}^{1/2}|t|^\kappa
       \left(\int_{\mathbb R}\ell_q(t+\sigma z)^2\,d\Phi_1(z)\right)\,dt
=\frac{c_{\mathrm{hi}}}{16}V_q<\infty.
\]
Tonelli's theorem and the upper design envelope therefore give, initially
as nonnegative extended integrals,
\[
\begin{aligned}
E_P[\mathbf1\{X=x\}\ell_q(V)^2]
&=p_x\int_{-1/2}^{1/2}g_x(t)
       \left(\int_{\mathbb R}\ell_q(t+\sigma z)^2\,d\Phi_1(z)\right)\,dt\\
&\le c_{\mathrm{hi}}\int_{-1/2}^{1/2}|t|^\kappa
       \left(\int_{\mathbb R}\ell_q(t+\sigma z)^2\,d\Phi_1(z)\right)\,dt\\
&=V_q(K)<\infty,
\end{aligned}
\]
where we used \(p_x\le1\).
By \cref{obj:lem:realized-dose-mean}, \(0\le Y\le1\) almost surely, so
\[
E_P[\mathbf1\{X=x\}Y^2\ell_q(V)^2]
\le E_P[\mathbf1\{X=x\}\ell_q(V)^2]
\le V_q(K).
\]
Both observable marks are thus square integrable. For the frequency mark,
\[
0\le\mathbf1\{X=x\}\le1,
\qquad
\operatorname{Var}_P(\mathbf1\{X=x\})\le\frac14,
\]
by the variance bound for a random variable in \([0,1]\).

Linearity of expectation gives
\[
E_{Q_P^n}\widehat D_x(q)=D_x(q),\qquad
E_{Q_P^n}\widehat M_x(q)=M_x(q),\qquad
E_{Q_P^n}\widehat p_x=p_x.
\]
Independence of the sample records makes the mixed covariance terms vanish.
Hence
\[
\begin{aligned}
\operatorname{Var}_{Q_P^n}(\widehat D_x(q))
&=\frac1{N_1^2}\sum_{i\in\mathcal I_1}
  \operatorname{Var}_P(\mathbf1\{X=x\}\ell_q(V))
\le\frac{V_q(K)}{N_1},\\
\operatorname{Var}_{Q_P^n}(\widehat M_x(q))
&=\frac1{N_2^2}\sum_{i\in\mathcal I_2}
  \operatorname{Var}_P(\mathbf1\{X=x\}Y\ell_q(V))
\le\frac{V_q(K)}{N_2},\\
\operatorname{Var}_{Q_P^n}(\widehat p_x)
&=\frac1{N_0^2}\sum_{i\in\mathcal I_0}
  \operatorname{Var}_P(\mathbf1\{X=x\})
\le\frac1{4N_0}.
\end{aligned}
\]
Here the first two inequalities also use
\(\operatorname{Var}_P(\cdot)\le E_P[(\cdot)^2]\).
These averages and their absolute centered errors are square integrable,
and therefore integrable.
% lean: CausalSmith.Stat.NoisydoseWeakdesignTransition.Dhat_population_moments
% lean: CausalSmith.Stat.NoisydoseWeakdesignTransition.Mhat_population_moments
% lean: CausalSmith.Stat.NoisydoseWeakdesignTransition.phat_population_moments

\item
Distinct modulo-three blocks are disjoint: an index in both
\(\mathcal I_0\) and \(\mathcal I_r\) would have remainder both \(0\) and
\(r\). Under the product experiment, the records indexed by disjoint
blocks are independent. Taking measurable functions of those records
preserves independence, so, separately for the denominator and numerator
blocks,
\[
\widehat p_x\perp\Delta_{D,x},
\qquad
\widehat p_x\perp\Delta_{M,x}
\qquad\text{under }Q_P^n.
\]
The preceding integrability and independence give integrable products and
the factorizations
\[
E_{Q_P^n}[\widehat p_x\Delta_{D,x}]
=p_xE_{Q_P^n}\Delta_{D,x},
\qquad
E_{Q_P^n}[\widehat p_x\Delta_{M,x}]
=p_xE_{Q_P^n}\Delta_{M,x}.
\]
% lean: CausalSmith.Stat.NoisydoseWeakdesignTransition.phat_block_error_factorization

\item
Nonnegativity of the variance of an absolute centered error bounds its
mean by the standard deviation of the original average. For example,
unbiasedness gives
\[
\begin{aligned}
0
&\le\operatorname{Var}_{Q_P^n}(\Delta_{D,x})\\
&=E_{Q_P^n}\Delta_{D,x}^2
  -(E_{Q_P^n}\Delta_{D,x})^2\\
&=\operatorname{Var}_{Q_P^n}(\widehat D_x(q))
  -(E_{Q_P^n}\Delta_{D,x})^2.
\end{aligned}
\]
Since the expected absolute error is nonnegative, taking square roots
yields the first of the following bounds; the same argument gives the
other two:
\[
\begin{aligned}
E_{Q_P^n}\Delta_{D,x}
&\le\sqrt{\operatorname{Var}_{Q_P^n}(\widehat D_x(q))}
\le\sqrt{\frac{V_q(K)}{N_1}},\\
E_{Q_P^n}\Delta_{M,x}
&\le\sqrt{\operatorname{Var}_{Q_P^n}(\widehat M_x(q))}
\le\sqrt{\frac{V_q(K)}{N_2}},\\
E_{Q_P^n}\Delta_{p,x}
&\le\sqrt{\operatorname{Var}_{Q_P^n}(\widehat p_x)}
\le\sqrt{\frac1{4N_0}}.
\end{aligned}
\]
Because \(n\le5N_r\), \(n>0\), and \(V_q(K)\ge0\),
\[
\frac{V_q(K)}{N_r}\le9\frac{V_q(K)}{n},
\qquad
\frac1{4N_r}\le\frac3n.
\]
Monotonicity of the square root consequently gives
\[
E_{Q_P^n}\Delta_{D,x}\le3\sqrt{\frac{V_q(K)}{n}},
\qquad
E_{Q_P^n}\Delta_{M,x}\le3\sqrt{\frac{V_q(K)}{n}},
\qquad
E_{Q_P^n}\Delta_{p,x}\le\sqrt{\frac3n}.
\]
These inequalities also hold when \(V_q(K)=0\).
% lean: CausalSmith.Stat.NoisydoseWeakdesignTransition.expected_abs_deviation_le_sqrt_variance
% lean: CausalSmith.Stat.NoisydoseWeakdesignTransition.inverse_moments_expected_abs_error
% lean: CausalSmith.Stat.NoisydoseWeakdesignTransition.split_standard_deviation_le

\item
Partitioning the bounded potential outcome \(Y(a_0)\) by the two strata
and using \cref{obj:def:conditional-potential-mean} gives
\[
\theta(P)
=\sum_{x\in\mathcal X}E_P[\mathbf1\{X=x\}Y(a_0)]
=\sum_{x\in\mathcal X}p_x\mu_x(a_0).
\]
The estimator therefore satisfies the exact decomposition
\[
\widehat\theta_q-\theta(P)
=\sum_{x\in\mathcal X}\widehat p_x
       \bigl(\widehat\mu_x(q)-\mu_x(a_0)\bigr)
 +\sum_{x\in\mathcal X}(\widehat p_x-p_x)\mu_x(a_0).
\]
The frequencies are nonnegative by their defining averages. Applying the
triangle inequality gives, for every observed sample,
\[
|\widehat\theta_q-\theta(P)|
\le
\sum_{x\in\mathcal X}\widehat p_x
       |\widehat\mu_x(q)-\mu_x(a_0)|
+\left|\sum_{x\in\mathcal X}(\widehat p_x-p_x)\mu_x(a_0)\right|.
\]
% lean: CausalSmith.Stat.NoisydoseWeakdesignTransition.weightEstimator_error_decomposition

\item
For every stratum, \cref{obj:lem:positive-ratio} supplies
\[
D_x(q)\ge b_q(K)>0,\qquad
m_{x,q}:=\frac{M_x(q)}{D_x(q)},\qquad
|m_{x,q}-\mu_x(a_0)|\le B_q(K).
\]
We also need the range of \(m_{x,q}\). The conditional-mean identity in
\cref{obj:lem:realized-dose-mean}, applied by the conditional-expectation
pull-out rule to the integrable inverse-weight mark, gives
\[
M_x(q)
=E_P[\mathbf1\{X=x\}\mu_X(A)\ell_q(V)].
\]
Indeed, law integrability and the bounded outcome and mean make both
products integrable. Conditioning on the positive-mass stratum, using
the Gaussian product law, and applying Fubini and pointwise inversion
now yield
\[
D_x(q)
=p_x\int_{-1/2}^{1/2}g_x(t)
       \left(\int_{\mathbb R}\ell_q(t+\sigma z)\,d\Phi_1(z)\right)\,dt
=p_x\int_{-1/2}^{1/2}g_x(t)q(t)\,dt,
\]
and
\[
\begin{aligned}
M_x(q)
&=p_x\int_{-1/2}^{1/2}g_x(t)\mu_x(a_0+t)
       \left(\int_{\mathbb R}\ell_q(t+\sigma z)\,d\Phi_1(z)\right)\,dt\\
&=p_x\int_{-1/2}^{1/2}g_x(t)q(t)\mu_x(a_0+t)\,dt.
\end{aligned}
\]
Absolute integrability under \(P\) persists under each normalized
stratum restriction, which justifies these uses of Fubini. The
conditional mean is measurable on the dose interval by its Hölder
condition.

The design envelopes and nonnegativity of \(q\) give, almost everywhere
on \([-1/2,1/2]\),
\[
0\le g_x(t)q(t)\le c_{\mathrm{hi}}|t|^\kappa q(t),
\qquad
\int_{-1/2}^{1/2}g_x(t)q(t)\,dt
\ge c_{\mathrm{lo}}I_\kappa(q)>0.
\]
Thus this weight is integrable and has positive mass. Since
\(a_0+t\in[0,1]\), \cref{obj:ass:mean-range} gives
\[
\begin{aligned}
0
&\le\int_{-1/2}^{1/2}g_x(t)q(t)\mu_x(a_0+t)\,dt\\
&\le\int_{-1/2}^{1/2}g_x(t)q(t)\,dt.
\end{aligned}
\]
Cancelling \(p_x>0\) in the population ratio and dividing by the positive
weight mass proves
\[
m_{x,q}
=\frac{\displaystyle\int_{-1/2}^{1/2}g_x(t)q(t)\mu_x(a_0+t)\,dt}
       {\displaystyle\int_{-1/2}^{1/2}g_x(t)q(t)\,dt}
\in[0,1].
\]

For every observed sample, define the floored denominator
\[
\widehat D_x^{\mathrm{floor}}(q)
:=\max\{\widehat D_x(q),b_q(K)/2\}.
\]
It satisfies
\[
\widehat D_x^{\mathrm{floor}}(q)\ge b_q(K)/2>0,
\qquad
|\widehat D_x^{\mathrm{floor}}(q)-D_x(q)|
\le|\widehat D_x(q)-D_x(q)|=\Delta_{D,x}.
\]
The latter follows because \(D_x(q)\ge b_q(K)\), and taking the maximum
with \(b_q(K)/2\) contracts distance to any number at least \(b_q(K)/2\).

The clipping construction in \cref{obj:def:total-estimator} is
\[
\widehat\mu_x(q)
=\Pi_{[0,1]}
 \left(\frac{\widehat M_x(q)}
 {\widehat D_x^{\mathrm{floor}}(q)}\right),
\qquad
\Pi_{[0,1]}(u)=\max\{0,\min\{1,u\}\}.
\]
Projection onto this interval fixes \(m_{x,q}\) and contracts distance
to it. Moreover,
\[
\frac{\widehat M_x(q)}{\widehat D_x^{\mathrm{floor}}(q)}-m_{x,q}
=
\frac{\widehat M_x(q)-M_x(q)
      +m_{x,q}\bigl(D_x(q)-\widehat D_x^{\mathrm{floor}}(q)\bigr)}
     {\widehat D_x^{\mathrm{floor}}(q)}.
\]
Using \(|m_{x,q}|\le1\), the floor bound, and the triangle inequality
therefore gives
\[
|\widehat\mu_x(q)-m_{x,q}|
\le\frac{\Delta_{M,x}+\Delta_{D,x}}
          {\widehat D_x^{\mathrm{floor}}(q)}
\le\frac2{b_q(K)}(\Delta_{M,x}+\Delta_{D,x}).
\]
Combining this with the localization bound yields
\[
|\widehat\mu_x(q)-\mu_x(a_0)|
\le B_q(K)+\frac2{b_q(K)}(\Delta_{M,x}+\Delta_{D,x}).
\]
Every denominator used here is positive for every observed sample.
% lean: CausalSmith.Stat.NoisydoseWeakdesignTransition.localizedRatio_mem_Icc
% lean: CausalSmith.Stat.NoisydoseWeakdesignTransition.muhat_error_le

\item
Each record belongs to exactly one of the two strata, so
\[
\widehat p_x\ge0,
\qquad
\sum_{x\in\mathcal X}\widehat p_x
=\frac1{N_0}\sum_{i\in\mathcal I_0}
       \sum_{x\in\mathcal X}\mathbf1\{X_i=x\}
=1.
\]
Multiplying the stratum error bounds by these frequencies and summing
gives
\[
\sum_{x\in\mathcal X}\widehat p_x
       |\widehat\mu_x(q)-\mu_x(a_0)|
\le
B_q(K)+\frac2{b_q(K)}\sum_{x\in\mathcal X}
       \bigl(\widehat p_x\Delta_{M,x}
             +\widehat p_x\Delta_{D,x}\bigr).
\]
Also, \cref{obj:ass:mean-range} implies
\(|\mu_x(a_0)|\le1\), whence
\[
\left|\sum_{x\in\mathcal X}(\widehat p_x-p_x)\mu_x(a_0)\right|
\le\sum_{x\in\mathcal X}|\widehat p_x-p_x|\,|\mu_x(a_0)|
\le\sum_{x\in\mathcal X}\Delta_{p,x}.
\]
Substitution into the error decomposition proves the pointwise bound
\[
|\widehat\theta_q-\theta(P)|
\le
B_q(K)+\frac2{b_q(K)}\sum_{x\in\mathcal X}
       \bigl(\widehat p_x\Delta_{M,x}
             +\widehat p_x\Delta_{D,x}\bigr)
+\sum_{x\in\mathcal X}\Delta_{p,x}.
\]
The right side is integrable by the moment bounds and the block
factorizations. Integrating this bound and splitting the finite sums
therefore gives
\[
\begin{aligned}
E_{Q_P^n}|\widehat\theta_q-\theta(P)|
\le B_q(K)
&+\frac2{b_q(K)}\sum_{x\in\mathcal X}
 E_{Q_P^n}\bigl[\widehat p_x\Delta_{M,x}
                  +\widehat p_x\Delta_{D,x}\bigr]\\
&+\sum_{x\in\mathcal X}E_{Q_P^n}\Delta_{p,x}.
\end{aligned}
\]
% lean: CausalSmith.Stat.NoisydoseWeakdesignTransition.phat_sum_eq_one
% lean: CausalSmith.Stat.NoisydoseWeakdesignTransition.public_error_certificate

\item
The block factorizations and deviation bounds imply, for each stratum,
\[
\begin{aligned}
E_{Q_P^n}\bigl[\widehat p_x\Delta_{M,x}
                 +\widehat p_x\Delta_{D,x}\bigr]
&=p_x\bigl(E_{Q_P^n}\Delta_{M,x}
           +E_{Q_P^n}\Delta_{D,x}\bigr)\\
&\le p_x\,6\sqrt{\frac{V_q(K)}{n}}.
\end{aligned}
\]
Unbiasedness of the frequencies and their samplewise normalization give
\[
\sum_{x\in\mathcal X}p_x
=\sum_{x\in\mathcal X}E_{Q_P^n}\widehat p_x
=E_{Q_P^n}\!\left[\sum_{x\in\mathcal X}\widehat p_x\right]
=1.
\]
Consequently,
\[
\sum_{x\in\mathcal X}
 E_{Q_P^n}\bigl[\widehat p_x\Delta_{M,x}
                  +\widehat p_x\Delta_{D,x}\bigr]
\le6\sqrt{\frac{V_q(K)}{n}}.
\]
There are two strata, and hence
\[
\sum_{x\in\mathcal X}E_{Q_P^n}\Delta_{p,x}
\le2\sqrt{\frac3n}.
\]
Finally, \(2/b_q(K)\ge0\), so the integrated pointwise bound yields
\[
\begin{aligned}
E_{Q_P^n}|\widehat\theta_q-\theta(P)|
&\le B_q(K)+\frac2{b_q(K)}\,6\sqrt{\frac{V_q(K)}{n}}
             +2\sqrt{\frac3n}\\
&=B_q(K)+\frac{12}{b_q(K)}\sqrt{\frac{V_q(K)}{n}}
             +2\sqrt{\frac3n}
=A_q(K).
\end{aligned}
\]
% lean: CausalSmith.Stat.NoisydoseWeakdesignTransition.public_error_certificate
\end{enumerate}
\end{proof}
